\question[5]
Ein Ordner kann in einem anderen Ordner liegen.
\begin{lstlisting}
class Ordner:
    def __init__(self, name, ober=None):
        self.name = name
        self.ober = ober

c = Ordner('C:')
s = Ordner('Schule', c)
i = Ordner('Informatik', s)
\end{lstlisting}
\begin{parts}
\part Zeichne das Klassendiagramm.
\part Ergänze die Methode \texttt{tiefe()}, die zurückgibt, in wie vielen Ordnern ein Ordner liegt.
\texttt{i.tiefe()} ergibt 2, \texttt{c.tiefe()} ergibt 0.
\end{parts}
\begin{solutionbox}{7cm}

\includegraphics[height=2.5cm]{\pfad/assoziation_02.png}

\begin{lstlisting}
    def tiefe(self):
        n = 0
        o = self.ober
        while o is not None:
            n += 1
            o = o.ober
        return n
\end{lstlisting}
\end{solutionbox}
